\chapter{Calibrated currents and mass comparison}\label{sec:calibrated-currents}
The smooth comparison of Proposition \ref{prop:calibration-minimizes} extends to currents and gives useful equality and convergence statements. Here $U\subset\mathbb R^N$ is open, $1\leq m\leq N$, and $\omega$ is a smooth closed $m$-form with comass at most one. The motivating source is \cite{ChMin}, Section 16, p.~46; current evaluation and mass use the conventions of \cite{Sim18}, Chapter 6, and Definition \ref{def:mass-current}.

\begin{definition}[A current calibrated by a form]\label{def:calibrated-current}
A compactly supported real $m$-current $T$ of finite mass in $U$ is calibrated by $\omega$ if
\[
 T(\omega)=\mathbf M(T).
\]
Evaluation uses a smooth cutoff equal to one near the support of $T$. For a locally integer rectifiable current, the local calibration condition is $\langle\omega(x),\tau_T(x)\rangle=1$ for $\|T\|$-almost every $x$. No global finite mass is required for this local condition. Compare Definition \ref{def:calibration} and \cite{ChMin}, Theorem 16.2.
\end{definition}

\begin{lemma}[The calibration defect and equality condition]\label{lem:calibration-defect}
For a compactly supported integer rectifiable current $T$ of finite mass,
\[
 \mathbf M(T)-T(\omega)
 =\int\bigl(1-\langle\omega,\tau_T\rangle\bigr)\,d\|T\|\geq0.
\]
Equality holds exactly when the local calibration condition holds almost everywhere.
\end{lemma}
\begin{proof}
Use the oriented rectifiable representation and mass formula in Definition \ref{def:integral-current}. Its unit tangent vector is simple, so the comass bound gives $\langle\omega,\tau_T\rangle\leq1$. The difference is the integral of a nonnegative function, and it vanishes if and only if that function vanishes almost everywhere. The sign of the orientation is part of the condition.
\end{proof}

\begin{theorem}[Calibrated currents minimize in their compact homology class]\label{thm:calibrated-current-minimizes}
Let $T$ be a compactly supported finite-mass real $m$-current calibrated by $\omega$ in $U$. If $S$ is another compactly supported finite-mass real $m$-current and $S-T=\partial R$ for a compactly supported $(m+1)$-current $R$ in $U$, then
\[
 \mathbf M(S)-\mathbf M(T)=\mathbf M(S)-S(\omega)\geq0.
\]
If $S$ is integer rectifiable, equality holds precisely when $S$ is also calibrated. This is the current form of the argument in \cite{ChMin}, Theorem 16.2, p.~46.
\end{theorem}
\begin{proof}
Choose a cutoff $0\leq\chi\leq1$ equal to one on a neighborhood of all three compact supports. Since $d\omega=0$,
\[
 S(\omega)-T(\omega)=\partial R(\chi\omega)
 =R(d\chi\wedge\omega+\chi d\omega)=0.
\]
The mass definition gives $S(\omega)\leq\mathbf M(S)$, since the cutoff form has comass at most one. Substitute $T(\omega)=\mathbf M(T)$; Lemma \ref{lem:calibration-defect} gives the equality statement. The filling is required to stay inside $U$.
\end{proof}

\begin{corollary}[Fixed-boundary minimization in Euclidean space]\label{cor:calibrated-euclidean-boundary}
Suppose $\omega$ is defined on all of $\mathbb R^N$. A compactly supported calibrated finite-mass $m$-current $T$, $m\geq1$, minimizes mass among all compactly supported finite-mass real $m$-currents $S$ with $\partial S=\partial T$.
\end{corollary}
\begin{proof}
The difference $Z=S-T$ is a compactly supported normal cycle: its mass is finite and its boundary is zero. Proposition \ref{prop:current-cone-bound} gives a compactly supported filling $R=C_aZ$ with $\partial R=Z$. Apply Theorem \ref{thm:calibrated-current-minimizes}. If $S,T$ are integral, this filling is integral too. No corresponding filling inside an arbitrary nonconvex open subset is inferred from this argument.
\end{proof}

\begin{proposition}[Weak limits of calibrated currents retain mass]\label{prop:calibrated-weak-limits}
Let $T_j$ be calibrated finite-mass real $m$-currents with supports in a common compact $K\subset U$, and suppose $T_j\rightharpoonup T$. Then $T$ is calibrated, is supported in $K$, and
\[
 \mathbf M(T_j)\longrightarrow\mathbf M(T).
\]
This conclusion does not by itself prove integrality of $T$.
\end{proposition}
\begin{proof}
A fixed cutoff equal to one near $K$ gives
$\mathbf M(T_j)=T_j(\omega)\to T(\omega)$, a finite limit. Tests outside $K$ show $\operatorname{spt}T\subset K$. Mass lower semicontinuity and the comass inequality give
\[
 \mathbf M(T)\leq\liminf_j\mathbf M(T_j)=T(\omega)\leq\mathbf M(T).
\]
Thus all terms agree. The finite mass of $T$ follows from the first inequality; the last evaluation is justified by that finite mass. Integrality requires an additional closure theorem, such as Theorem \ref{thm:integral-compactness} with its boundary-mass hypothesis.
\end{proof}

\begin{proposition}[Stationarity away from the boundary of a calibrated current]\label{prop:calibrated-current-stationary}
Let $T$ be a compactly supported calibrated integral $m$-current in $U$, and let $V_T$ be its rectifiable varifold obtained by forgetting orientation and retaining the positive integer multiplicity. Then $V_T$ is stationary in $U\setminus\operatorname{spt}\partial T$.
\end{proposition}
\begin{proof}
Take a smooth vector field $X$ compactly supported away from $\operatorname{spt}\partial T$, with flow $F_t$. For small $t$, the trace of this flow on the compact support stays inside $U$. The homotopy formula, Proposition \ref{prop:current-homotopy}, shows that $(F_t)_\#T-T$ is a compactly supported boundary in $U$. The term involving $\partial T$ vanishes because the flow is constant on a neighborhood of that boundary, so its derivative in the time direction is zero there. Theorem \ref{thm:calibrated-current-minimizes} gives $\mathbf M((F_t)_\#T)\geq\mathbf M(T)$ for both signs of $t$.

A diffeomorphism causes no overlap cancellation; the area formula expresses this mass as $\int J_m(F_t|S)\,dV_T(x,S)$. Differentiating at zero gives $\int\operatorname{div}_S X\,dV_T=0$. Bounded derivatives of the flow on a compact set justify differentiation under the finite measure. Smooth fields are dense in compactly supported $C^1$ fields, giving stationarity in the stated domain.
\end{proof}

\begin{example}[Why the homology condition matters]\label{ex:calibration-homology-condition}
In the flat torus $M=(\mathbb R/\mathbb Z)^{m+1}$, the oriented slice $\Sigma=(\mathbb R/\mathbb Z)^m\times\{0\}$ is calibrated by $\omega=dx_1\wedge\cdots\wedge dx_m$. It has empty boundary and positive area, while the zero chain has the same boundary. These two cycles are not homologous.
\end{example}
\begin{proof}
The angular one-forms are globally defined and closed; their wedge has comass one and restricts to the positive volume form of the slice. Hence $\int_\Sigma\omega=1$. If $\Sigma$ were a boundary, Stokes' theorem and $d\omega=0$ would force that integral to be zero. Thus the calibration theorem does not compare this cycle to zero in its homology class. In contrast, Corollary \ref{cor:calibrated-euclidean-boundary} uses a filling available in Euclidean space.
\end{proof}
